\section{Linnik's problem and where zeros enter}\label{gl:sec:linnik}

\subsection{The theorem and the records}

Linnik \cite{Linnik51,Linnik53} proved that there is an absolute constant $K$ such that every sufficiently large even integer $N$ can be written
\begin{equation}\label{gl:eq:linnik}
N = p_1 + p_2 + 2^{k_1} + \cdots + 2^{k_K},\qquad p_1,p_2\ \text{prime},\quad k_i\in\mathbb Z_{\ge0}.
\end{equation}
The 1951 paper assumes GRH and the 1953 paper does not; neither gives a value of $K$. Gallagher \cite{Gallagher75} recast the argument so that explicit constants became possible, and the values of $K$ that followed were lowered in a series of papers \cite{LLW98,LLW99,Li00,Wang99,Li01}, from $54000$ (Liu, Liu and Wang) to $1906$ (Li), and under GRH to $200$ (Liu, Liu and Wang), until Heath-Brown and Puchta \cite{HBP02} reached the values in the first line of the table below.

% CHECK: the rows of the table of records that the first version had before review-A M9 (Liu, Liu and Wang, Li, Wang, Li) are replaced by the values that HBP02 recall in their Section 1: 54000 by Liu, Liu and Wang (their reference 12, which is part II, Science in China Series A 41 (1998), 1255 to 1271, and is not in the bibliography, LLW98 being their part I), then 25000 (Li), 2250 (Wang) and 1906 (Li), and 200 under GRH (Liu, Liu and Wang, LLW99); the authors may add part II and restore the rows; Liu and Lu 12 unconditionally (LiuLu11) stays out, since its title concerns two squares of primes, and its bibliography entry is removed because nothing cites it (it is in misc/reviews/nick-linnik/bib-new.tex; question Q26 of misc/catalog/QUESTIONS.md); Pintz announced 12 unconditionally and 10 under GRH around 2000, and Elsholtz obtained 12 in unpublished work, both recorded in Section 1 of HBP02 (arXiv math/0201299v2, checked on 17 October 2026).
% CHECK: Platt and Trudgian (PlattTrudgian15): the source says that sharpening the constants they came within 0.1 of K=11 unconditionally and of K=6 under GRH; there is no pinpoint and the claim is not printed (review-A M9g).
The best values on record are the following.
\begin{center}
\begin{tabular}{lcc}
\toprule
Authors & $K$ unconditionally & $K$ under GRH\\
\midrule
Heath-Brown and Puchta \cite{HBP02} & $13$ & $7$\\
Pintz and Ruzsa \cite{PR03} & & $7$\\
Pintz and Ruzsa \cite{PR20} & $8$ &\\
\bottomrule
\end{tabular}
\end{center}
Platt and Trudgian \cite{PlattTrudgian15} sharpened the numerical constants in the method of \cite{HBP02}.

\subsection{Where \texorpdfstring{$K$}{K} is determined}\label{gl:sec:where}\label{gl:sec:oct}
% gl:sec:oct names the paragraph that closes this subsection; a paragraph has no number, so the label sits on the subsection heading and Section~\ref{gl:sec:oct} prints the number of the subsection.

Fix a small constant $\varpi>0$, write $e(x)=\exp(2\pi ix)$ and $L=\lfloor\log_2(N/2K)\rfloor$, and put
\[
\mathcal S(\alpha)=\sum_{\varpi N<p\le N}e(p\alpha),\qquad \mathcal D(\alpha)=\sum_{1\le\nu\le L}e(2^\nu\alpha).
\]
The number of representations \eqref{gl:eq:linnik} with both primes in $(\varpi N,N]$ and the exponents restricted to $1\le k_i\le L$, so that the powers of two add up to at most $N/2$, is
\begin{equation}\label{gl:eq:circle}
R(N)=\int_0^1 \mathcal S(\alpha)^2\,\mathcal D(\alpha)^K\,e(-N\alpha)\,d\alpha .
\end{equation}
Let $\mathfrak M\subset[0,1]$ be the set of major arcs, the neighbourhoods of rationals with small denominator, and $\mathfrak m=[0,1]\setminus\mathfrak M$ the minor arcs.

% CHECK: HBP02 state Lemma 7, the major arc lemma of step (1), for K at least 9 (equation (25) and the lemma), while their Theorem 2 has K=7 under GRH; Section 6 compares the minor arcs with the same constant 2.7895 for both values of K and does not say how the lemma is used for K=7, although the computation for the divisors 3 and 5 holds for K at least 7; the authors should confirm that this is intended. Checked against arXiv math/0201299v2 on 17 October 2026 (misc/reviews/nick-linnik/hbp02-check.md): the definitions of S, T and L, the constants 2.7895 (equation (25) and Lemma 7) and 13.968 (Lemmas 8 to 10), the exponents 263/308 (Theorem 3.1 of Vaughan's book, with P = N^(45/154 - 4 varpi)) and 3/4 (Lemma 12 of Baker and Harman, under GRH), E(0.863665) > 109/154 and E(0.722428) > 1/2 (Lemma 1, proved in Section 7), and the comparison 13.968 lambda^(K-2) < 2.7895 at the end of Section 6, where the paper prints K > 12.991 and K > 6.995 and the comparison gives 2 + 10.99 and 2 + 4.95.
We follow the argument of Heath-Brown and Puchta \cite{HBP02}, whose constants we quote with their pinpoints; the proof has five steps.
\begin{enumerate}
\item[(1)] \emph{Major arcs.} Gallagher's method, with the major arcs $q\le P$, $|\alpha-a/q|\le P/(qN)$ and $P\le N^{45/154-4\varpi}$, gives on $\mathfrak M$ a main term of at least
\[
2.7895\,(1-3\varpi)\,C_0\,N(\log2)^{-2}L^{K-2}
\]
(Lemma~7 of \cite{HBP02}, stated there for $K\ge9$), where $C_0=\prod_{p>2}\big(1-(p-1)^{-2}\big)=0.6601\ldots$ is the twin prime constant, and $2.7895$ is $2\big(1+\tfrac13(1-2^{-6})+\tfrac1{15}(1-4^{-6})\big)=2.78955\ldots$ rounded down, the contribution of the divisors $1$, $3$ and $5$ to the singular series of Goldbach's problem averaged over the powers of two (equation~(25) there). No hypothesis on zeros is used, and no explicit numerical zero free region. \cite{HBP02} control the error terms, character sums over primes, by the zero density estimates of Huxley and Jutila, $\sum_{r\le R}N(\sigma,r,T)\ll(R^2T)^{\kappa(\sigma)(1-\sigma)}$ with $\kappa(\sigma)=\tfrac3{2-\sigma}+\varpi$ for $\tfrac12\le\sigma\le\tfrac34$ and $\tfrac{12}5+\varpi$ for $\tfrac34\le\sigma\le1$, which also fix the restriction on $P$, and, for the few characters with a zero near $1$ and for the principal character, by the zero free regions of Iwaniec (for nonreal characters, with an implied constant that is not explicit) and of Vinogradov. Zeros close to $\sigma=1$ play no role.
\item[(2)] \emph{Minor arc mean square.} $\int_0^1|\mathcal S(\alpha)\mathcal D(\alpha)|^2d\alpha$ counts the solutions of $p_1-p_2=2^\mu-2^\nu$ with $\mu,\nu\le L$. Chen's sieve bound for the number of prime differences, $r(n)\le C_0C_1h(n)N(\log N)^{-2}$ with $C_1=7.8342$ (Theorem~3 of Chen's 1978 paper on Goldbach's problem and the sieve methods, as used in \cite{HBP02}), bounds the integral over $[0,1]$ from above, and a lower bound for the part over $\mathfrak M$ is subtracted, which gives $\int_{\mathfrak m}|\mathcal S\mathcal D|^2d\alpha\le(13.968+o(1))\,C_0N(\log2)^{-2}$ (Lemmas~8 to~10 of \cite{HBP02}).
\item[(3)] \emph{Minor arc sup bound for the primes.} $\mathcal S(\alpha)\ll N^{\theta+o(1)}$ on $\mathfrak m$, with $\theta=263/308$ unconditionally, from Theorem~3.1 of Vaughan's book on the Hardy Littlewood method with $P=N^{45/154-4\varpi}$, and $\theta=3/4$ under GRH, from Lemma~12 of Baker and Harman (Section~6 of \cite{HBP02}).
\item[(4)] \emph{Large deviations for the powers of two.} The set $\mathcal A_\lambda=\{\alpha:|\mathcal D(\alpha)|\ge\lambda L\}$ has measure $\ll N^{-E(\lambda)}$, with $E(0.863665)>\tfrac{109}{154}$ and $E(0.722428)>\tfrac12$ (Lemma~1 of \cite{HBP02}, proved there by a martingale argument suggested by Keith Ball). The contribution of $\mathfrak m\cap\mathcal A_\lambda$ is $\ll N^{-E(\lambda)+2\theta+o(1)}L^K$, negligible once $E(\lambda)>2\theta-1$, that is $\tfrac{109}{154}$ unconditionally and $\tfrac12$ under GRH.
\item[(5)] \emph{Comparison.} Off $\mathcal A_\lambda$, $|\mathcal D|<\lambda L$, so the contribution of $\mathfrak m\setminus\mathcal A_\lambda$ is at most $(\lambda L)^{K-2}\cdot13.968\,C_0N(\log2)^{-2}$, and the count is positive, for $N$ large and $\varpi$ small, when $13.968\,\lambda^{K-2}<2.7895$ (end of Section~6 of \cite{HBP02}). This needs $K-2>10.99$ at $\lambda=0.863665$ and $K-2>4.95$ at $\lambda=0.722428$, which gives $K=13$ unconditionally and $K=7$ under GRH.
\end{enumerate}
In this summary a hypothesis on zeros enters only at step~(3): GRH lowers $\theta$ to $\tfrac34$ and hence lowers $\lambda$. In the scheme of \cite{HBP02} the remaining gains come from the constants of steps~(1), (2) and (4), which do not depend on a hypothesis on zeros; Elsholtz is reported there to have obtained $K=12$ unconditionally, in unpublished work, by improving the constants of the first two steps. The refinements of Pintz and Ruzsa \cite{PR03,PR20} differ from this argument in a number of respects, as \cite{HBP02} remark, and we do not describe them. In the scheme of \cite{HBP02} an input that delivered GRH would be worth the passage from $K=13$ to $K=7$, and against the records it is the passage from $K=8$ to $K=7$.

\paragraph{The zero free half plane of October 2026.} In October 2026 OpenAI released a proof, produced with one of its AI models, that $\zeta(s)$ and every Dirichlet $L$ function have no zeros in $\Re s>\tfrac78$ \cite{OAI26a}, and a separate proof of the weaker half plane $\Re s>\tfrac{11}{12}$ \cite{OAI26b}. The claims have not been refereed (Section~\ref{sec:ch25:open}). A zero free region can act in step~(3), through the error term of the exponential sum, and in step~(1), through the zero density estimates. For step~(3), the exponent $\tfrac34$ under GRH comes from the lemma of Baker and Harman used in \cite{HBP02}: for $\alpha=a/q+\beta$ with $|\beta|\le N^{-1/2}$, GRH gives $\mathcal S(\alpha)\ll(\log N)\{N/q+N^{1/2}q^{1/2}+N(q|\beta|)^{1/2}\}$. By Dirichlet's theorem every $\alpha$ has such an approximation with $q\le N^{1/2}$ and $|\beta|\le1/(qN^{1/2})$, which makes each of the three terms at most $N^{3/4}$ once $q>N^{1/4}$, and $q\le N^{1/4}$ does not occur on the minor arcs because $P=N^{45/154-4\varpi}\ge N^{1/4}$. In the second term GRH enters as the factor $N^{1/2}$, the size of the error term in the explicit formula for $\psi(x,\chi)$. Replacing GRH by a zero free half plane $\Re s>\theta_0$ for all Dirichlet $L$ functions would replace that factor by $N^{\theta_0}$, and the same choice of approximation, with $q\le N^{1-\theta_0}$, would give $\theta=\tfrac12(1+\theta_0)$. This is a heuristic count: the lemma is proved under GRH and its analogue for a half plane is not proved here. It beats the unconditional exponent $\tfrac{263}{308}$ only if $\theta_0<\tfrac{109}{154}=0.7078$. The half plane $\Re s>\tfrac78$ gives $\tfrac{15}{16}$, worse than the unconditional exponent, so used this way it does not lower $K$. For step~(1), the constant $2.7895$ is arithmetic, and zeros close to $\sigma=1$ play no role in \cite{HBP02}, so excluding exceptional zeros is not what a half plane would add. Zeros enter the major arcs through the zero density estimates, which fix the restriction $P\le N^{45/154-4\varpi}$ by way of the bound $(\tfrac{45}{77}\kappa(\sigma)-1)(1-\sigma)\le\tfrac{31}{77}(1-\sigma)\le\tfrac{31}{154}$ for the exponent of the estimate (Section~3 there), taken at the left end $\sigma=\tfrac12$; a half plane $\Re s>\theta_0$ removes only the zeros to the right of $\theta_0$, so it does not relax this restriction. Whether a finer combination of the half plane with Vaughan's identity does better, on the major or the minor arcs, is Problem~\ref{gl:op:quasi}. In the Thorin language of this chapter, the half plane would give the Thorin property of the shifted clock at every basepoint $\alpha>\tfrac78$ for $\xi_F$ with $F$ abelian over $\mathbb Q$ (quadratic and cyclotomic fields in particular, by \eqref{gl:eq:cyclotomic}) and for every real character (Theorem~\ref{gl:thm:sinsq}; for $\xi$ this is the remark after Theorem~\ref{thm:poisson}). Liu \cite{Liu26b} claims the half plane $\Re s>B_\flat=0.87489$ for every Hecke $L$ function of finite order over every number field (Section~\ref{tx3:sec:signal}); if that unrefereed and unformalised claim stands, the same holds at every basepoint $\alpha>B_\flat$ for $\xi_F$ with $F$ any number field, abelian or not.
